%=============================================================
% Paper 9: Carrillo-Tudela & Visschers (2023), Econometrica
%=============================================================
\chapter{Carrillo-Tudela \& Visschers (2023)}

\begin{paperbox}{Unemployment and Endogenous Reallocation Over the Business
Cycle}
\textbf{Authors:} Carlos Carrillo-Tudela (Essex) \& Ludo Visschers
(Edinburgh/UC3M)\\[3pt]
\textbf{Journal:} \textit{Econometrica} 91(3): 1119--1153\\[3pt]
\textbf{Year:} 2023 \quad\textbf{Field:} Labour economics; Macroeconomics\\[3pt]
\textbf{Classification:} Structural (multisector business cycle model + SMM)
\end{paperbox}

%------------------------------------------------------------
\section{Research Question}

To what extent does the cyclicality of occupational mobility among the
unemployed shape the aggregate unemployment rate and its duration distribution?
Is cyclical unemployment driven primarily by occupation-wide productivity
differences (net reallocation) or by individual workers' changing career
prospects within occupations (idiosyncratic career shocks)?

%------------------------------------------------------------
\section{Motivation}

On average, 44\% of unemployed U.S.\ workers change major occupational group
(MOG) at reemployment. Occupational movers take longer to find jobs, and in
recessions their duration increases $\approx$35\% more than stayers'. Yet
standard multisector models (Lucas--Prescott 1974; Lilien 1982) emphasise
\emph{net} reallocation across sectors as the main driver of unemployment
fluctuations---and in such models, gross mobility is typically constant or
countercyclical. Carrillo-Tudela and Visschers document that gross occupational
mobility is \emph{procyclical} while net mobility is \emph{countercyclical},
and that the dominant component is \emph{excess mobility} (moves that cancel
out at the occupation level). This motivates a new theory where idiosyncratic
career shocks, not occupation-wide differences, drive cyclical unemployment.

%------------------------------------------------------------
\section{Main Contribution}

\begin{enumerate}
  \item \textbf{New empirical facts (SIPP, corrected for miscoding):}
    \begin{itemize}
      \item 44.5\% of unemployed workers change MOG at reemployment.
      \item Gross mobility is \emph{procyclical}; net mobility is
        \emph{countercyclical}.
      \item 90\% of gross mobility is excess mobility (procyclical).
      \item Gross mobility increases with unemployment duration (mobility-
        duration profile).
    \end{itemize}
  \item \textbf{Multisector business cycle model with idiosyncratic career
    shocks} ($z$): workers face both occupation-wide productivity differences
    ($p_o$) and idiosyncratic career shocks ($z$). A ``wait zone'' between the
    separation and reallocation cutoffs generates rest/wait unemployment.
  \item \textbf{Central quantitative finding:} Idiosyncratic career shocks
    ($z$), not occupation-wide differences ($p_o$), are the dominant driver
    of cyclical unemployment fluctuations.
\end{enumerate}

%------------------------------------------------------------
\section{Relationship to the Literature}

Challenges the Lucas--Prescott (1974)/Lilien (1982)/Rogerson (1987)
multisector framework. Related to Alvarez--Shimer (2011) rest/search model
(but extends to dynamic business cycle). Complements Pilossoph (2014) and
Chodorow-Reich--Wieland (2020)---which also find muted effects of net
reallocation---but those papers have countercyclical or acyclical gross
mobility, inconsistent with the data. Distinguishes from Kambourov--Manovskii
(2008) by focusing exclusively on workers who go through unemployment.

%------------------------------------------------------------
\section{Theoretical Framework and Economic Mechanism}

\subsection*{Environment}

Discrete time. Workers differ in: idiosyncratic productivity $z_t$ (Markov,
common across occupations---``career match'' or career prospect); occupation-
specific human capital $x_t$ (accumulates when employed via learning-by-doing;
depreciates during unemployment). Occupation-wide productivity $p_{o,t}$ and
aggregate productivity $A_t$ follow separate Markov processes.

\subsection*{Value Function---Unemployed Worker}

\begin{equation}
  W^U(\omega) = b + \beta E_{\omega'}\!\left[\max_{\rho(\omega')}\!
    \left\{\rho\,R(\omega') + (1-\rho)\!\left[\lambda(\theta(\omega'))W^E(\omega')
    + (1-\lambda(\theta(\omega')))W^U(\omega')\right]\right\}\right],
  \label{eq:ct_wu}
\end{equation}
where $\rho\in\{0,1\}$ is the reallocation decision, $R(\omega)$ is the
expected net value of searching across occupations, $\theta(\omega)$ is labour
market tightness in submarket $(z,x)$ of occupation $o$.

\subsection*{Separation and Reallocation Cutoffs}

Within each occupation, the model yields two productivity thresholds:
\begin{align}
  &\text{\textbf{Separation cutoff}} \quad z_s:
    &&W^E(\omega) < W^U(\omega) \;\Leftrightarrow\; z < z_s \\
  &\text{\textbf{Reallocation cutoff}} \quad z_r:
    &&R(\omega) > W^U(\omega) \;\Leftrightarrow\; z < z_r
\end{align}
The calibration finds $z_s > z_r$---creating a \textbf{wait zone}
$(z_r, z_s)$ where workers who separate prefer to remain attached to their
occupation rather than switch, because they expect their career prospect $z$
to recover.

\subsection*{Cyclical Widening of the Wait Zone}

In recessions, aggregate productivity $A$ falls $\Rightarrow$ wages fall
$\Rightarrow$ both $z_s$ and $z_r$ rise. Because $z_s$ rises more than $z_r$,
the gap $z_s - z_r$ \emph{widens}, pulling more workers into the wait zone
$\Rightarrow$ more rest/wait unemployment $\Rightarrow$ longer spells,
disproportionately for movers.

\subsection*{Imperfect Directed Search Across Occupations}

Workers allocate search effort $\{s_{\tilde{o}}\}$ across remaining
occupations with $\sum_{\tilde{o}\in O^-}s_{\tilde{o}}=1$. Concavity of the
matching function creates a trade-off between targeting the best occupation
and the total probability of receiving any offer. This generates
\emph{excess mobility}: workers switch occupations even without systematic
occupation-wide differences, driven by idiosyncratic career prospects.

%------------------------------------------------------------
\section{Data}

\begin{itemize}
  \item \textbf{SIPP:} 1984--2008 panels (covering 1983--2014); EUE
    (employment $\to$ unemployment $\to$ employment) spells; excludes
    self-employed, military, agriculture.
  \item \textbf{Miscoding correction:} Novel degarbling matrix $\Gamma$
    identified from the 1985$\to$1986 SIPP interviewing method change
    (independent to dependent). True occupational stayers have a
    $\approx$20\% chance of appearing as movers. Assumptions: (A1) independent
    classification errors; (A2) detailed balance in miscoding; (A3) strict
    diagonal dominance.
  \item \textbf{Cross-validation:} PSID 1968--1997; CPS 1979--2019.
  \item \textbf{Occupational classification:} 21 major groups (2000 SOC);
    also task-based (NRC, RC, NRM, RM).
  \item \textbf{Calibration:} Simulated Method of Moments (SMM).
\end{itemize}

%------------------------------------------------------------
\section{Empirical Strategy}

\begin{enumerate}
  \item \textbf{Classification error correction:} Degarble observed flows
    $M^I=\Gamma'M\Gamma$ to recover true flows $M=(\Gamma^{-1})'M^I\Gamma^{-1}$.
  \item \textbf{Mobility-duration profile:} For duration $x$, compute the
    proportion of workers who changed MOG among those with spells $\geq x$
    months (Figure~1).
  \item \textbf{Cyclicality regressions (Table~I):} Regress (log) gross
    mobility on HP-filtered unemployment rate; report across SIPP
    (corrected/uncorrected), CPS, with/without controls.
  \item \textbf{Net/excess decomposition:} Net mobility
    $= \tfrac{1}{2}\sum_i|E_{-i}UE_i - E_iUE_{-i}|/EUE$; excess
    $=$ gross $-$ net.
  \item \textbf{Model calibration:} SMM on moments including gross mobility
    rate, mobility-duration slope, cyclicality, unemployment volatility,
    duration distribution, Beveridge curve.
\end{enumerate}

%------------------------------------------------------------
\section{Identification Strategy}

The causal claim---that idiosyncratic career shocks ($z$), not occupation-wide
differences ($p_o$), drive cyclical unemployment---is identified through model
counterfactuals (shutting off $z$ vs.\ $p_o$ shocks). The empirical
identification rests on:
\begin{itemize}
  \item Gross flows order-of-magnitude larger than net flows $\Rightarrow$
    occupation-wide differences cannot account for the data.
  \item Procyclical gross mobility (not net) $\Rightarrow$ models with
    countercyclical net reallocation are inconsistent.
  \item SMM estimates identify the $z$ shock process from the
    mobility-duration profile slope.
\end{itemize}

The miscoding correction identification relies on the detailed balance
assumption (A2), tested and supported using SIPP--PSID--CPS data.

%------------------------------------------------------------
\section{Main Results}

\subsection*{Empirical Results}

\begin{center}
\small
\begin{tabular}{lc}
\toprule
\textbf{Statistic} & \textbf{Value} \\
\midrule
Gross occupational mobility at reemployment & 44.5\% \\
Mobility at $\geq$9 months of duration & 54.6\% \\
Average excess mobility (of gross) & 90\% \\
Average net mobility & 4.3\% \\
Cyclicality of gross mobility (SIPP, corrected): elast.\ & $-0.145$ \\
Duration gap (mover vs.\ stayer): expansion & 0.4 months \\
Duration gap: recession & 0.9 months \\
\bottomrule
\end{tabular}
\end{center}

Net mobility is \emph{countercyclical}; excess mobility is \emph{procyclical}
and drives the procyclicality of gross mobility.

\subsection*{Model Results}

\begin{itemize}
  \item Quantitatively matches: procyclical gross mobility, countercyclical
    net mobility, downward-sloping Beveridge curve, unemployment volatility,
    duration distribution, cyclicality of job-finding and separation rates.
  \item In recessions: rest/wait unemployment becomes relatively more
    prominent; search unemployment dominates in expansions.
  \item \textbf{Counterfactual:} Shutting off $z$ shocks eliminates most of
    the cyclical variation in unemployment; shutting off $p_o$ differences
    has little effect.
\end{itemize}

%------------------------------------------------------------
\section{Economic Magnitude and Interpretation}

The duration gap between movers and stayers grows from 0.4 to 0.9 months
from expansion to recession---a 40\% contribution to the overall increase in
average unemployment duration across the cycle. This means roughly two-fifths
of the aggregate duration increase is attributable to the worsening situation
of occupational movers specifically. The finding that 90\% of gross mobility
is excess (cancelling out at the occupation level) implies that the standard
narrative of ``structural unemployment'' driven by workers fleeing hard-hit
sectors accounts for only about 10\% of actual occupational transitions.

%------------------------------------------------------------
\section{Mechanisms}

The central mechanism is the \textbf{endogenous widening of the wait zone}
in recessions:
\begin{enumerate}
  \item $A\downarrow$ $\Rightarrow$ wages fall $\Rightarrow$ both
    $z_s$ and $z_r$ rise.
  \item $z_s$ rises more than $z_r$ (because option value of staying in the
    occupation is more affected by falling wages than the option value of
    reallocation, which provides a fresh $z$ draw).
  \item More workers have career prospects $z\in(z_r, z_s)$ $\Rightarrow$
    they separate but \emph{wait} rather than switch.
  \item Waiting motive is stronger for occupational movers (who must redraw
    $z$ upon switching) $\Rightarrow$ their spells lengthen more in recessions.
\end{enumerate}

The \textbf{excess mobility mechanism}: workers who draw low $z$ from the
ergodic distribution after separation find their occupation unattractive and
search in other occupations. Since most eventually draw a new $z$ (often
returning to prior occupation), these moves cancel at the occupation level
$\Rightarrow$ large gross but small net flows.

%------------------------------------------------------------
\section{Robustness Checks}

Extensive (Online Supplementary Material + Supplementary Appendices):
\begin{itemize}
  \item Results hold under alternative classifications (21-group SOC,
    task-based NRC/RC/NRM/RM).
  \item Procyclicality robust using SIPP (corrected/uncorrected), CPS
    1979--2019, with/without demographic controls.
  \item Mobility-duration profile robust to nonemployment spells, alternative
    spell definitions, CPS 1994-redesign restriction.
  \item Classification error assumptions (A1--A3) tested with PSID and CPS.
  \item Non-targeted moment: repeat mobility (63.4\% of stayers remain stayers;
    54.4\% of movers become movers again)---matched by the model.
\end{itemize}

%------------------------------------------------------------
\section{Limitations}

\begin{enumerate}
  \item \textbf{EUE spells only:} Workers who never return to employment
    are excluded; may miss extreme reallocation cases.
  \item \textbf{Coarse occupation classification:} 21 major groups may miss
    finer mobility patterns.
  \item \textbf{BRE restriction:} Block Recursive Equilibrium restricts value
    functions to individual states plus aggregate---a technical tractability
    assumption.
  \item \textbf{Exogenous wage determination:} Nash bargaining; richer wage
    dynamics (posting, rigidity) may alter quantitative results.
  \item \textbf{US-specific calibration:} The $z$ shock process and mobility
    rates are calibrated to U.S.\ data; generalisability unclear.
\end{enumerate}

%------------------------------------------------------------
\section{Policy Implications}

\begin{itemize}
  \item Idiosyncratic career shocks dominate $\Rightarrow$ standard
    \emph{sectoral retraining} policies address only $\approx$10\% of
    occupational transitions (net flows). Most mobility is driven by
    individual career dynamics.
  \item The wait zone mechanism implies workers appearing ``structurally
    unemployed'' are often rationally waiting for career recovery. Policies
    forcing rapid reallocation (e.g.\ strict UI time limits) may push workers
    to accept worse occupational matches prematurely.
  \item Occupational movers have 35\% longer duration increases in recessions
    $\Rightarrow$ supports differential UI policies for occupational switchers
    (longer or more generous benefits during transitions).
\end{itemize}

%------------------------------------------------------------
\section{Overall Assessment}

An exceptional paper combining careful empirical measurement (novel miscoding
correction), a tractable theoretical model, and rigorous SMM calibration.
The main empirical fact---procyclical gross mobility, countercyclical net
mobility, excess mobility dominates---overturns a widely held prior. The
model's success in matching this complex pattern simultaneously with the
cyclical properties of unemployment duration is impressive.

\textbf{Quality: Outstanding. A landmark contribution to occupational
mobility, unemployment dynamics, and business cycles.}

%------------------------------------------------------------
\section{Relevance to My Research}

Highly relevant for RDC work. Canadian LEED/LEAP data may allow measurement
of occupational mobility among the unemployed, testing whether the procyclical
gross/countercyclical net finding holds in Canada. The wait-zone mechanism
has direct relevance for resource-sector workers during commodity busts.

%------------------------------------------------------------
\section{Potential Research Extensions}

\begin{itemize}
  \item Test the mobility-duration profile and its cyclicality for Canada
    using LWF/LEED data.
  \item Study whether the wait zone is wider for resource-sector workers in
    Canada during commodity busts (high occupation-specific human capital
    $\Rightarrow$ stronger waiting motive).
\end{itemize}

%------------------------------------------------------------
\section{Research Gaps}

\begin{itemize}
  \item The paper does not model long-run secular changes in occupational
    mobility (automation, polarisation); the excess mobility mechanism might
    interact with technological change.
  \item Firm actions (R&D, capital investment) could endogenously affect
    workers' career prospects $z$---this feedback is absent.
\end{itemize}

%------------------------------------------------------------
\section{Methodological Lessons}

\begin{enumerate}
  \item The degarbling matrix $\Gamma$ identified from a data collection
    design change (SIPP interviewing method) is a powerful and broadly
    applicable tool for measuring occupational mobility free of coding error.
  \item Decomposing gross mobility into net and excess components is essential:
    the two have completely different cyclical properties; aggregating them
    would obscure the key empirical patterns.
  \item Block Recursive Equilibrium is an effective approach to tractability
    in rich multisector heterogeneous-agent models---it rules out the joint
    productivity distribution as a state variable while preserving full
    individual-level heterogeneity.
\end{enumerate}

%------------------------------------------------------------
\section*{Five-Sentence Summary}

Carrillo-Tudela and Visschers (2023) document, using SIPP data with a novel
miscoding correction, that 44\% of unemployed U.S.\ workers change major
occupational group at reemployment, gross occupational mobility among the
unemployed is procyclical, net mobility is countercyclical, and 90\% of gross
mobility is excess mobility---patterns inconsistent with standard multisector
models. They develop a multisector business cycle model with idiosyncratic
career shocks ($z$), occupation-specific human capital, and DMP matching
frictions; within each occupation the separation cutoff exceeds the reallocation
cutoff, creating a wait zone where separated workers prefer to remain attached
to their occupation. In recessions, aggregate productivity falls, the wait zone
widens endogenously, more workers enter rest/wait unemployment, and duration
spells lengthen disproportionately for occupational movers---generating the
observed 35\% larger duration increase for movers than stayers in downturns.
The model is calibrated by SMM and quantitatively matches the mobility-duration
profile, procyclical gross mobility, countercyclical net mobility, unemployment
volatility, and the Beveridge curve simultaneously. The core finding is that
cyclical unemployment is driven by workers' changing idiosyncratic career
prospects, not occupation-wide differences---implying standard sectoral
retraining policies address only a small fraction ($\approx$10\%) of actual
occupational transitions.

\vspace{6pt}
\begin{quickref}
\begin{tabularx}{\linewidth}{@{}lX@{}}
\textbf{Key contribution} &
  First model to simultaneously reproduce procyclical gross and countercyclical
  net occupational mobility among the unemployed; shows idiosyncratic career
  shocks (not occupation-wide differences) drive cyclical unemployment through
  an endogenous wait-zone mechanism. \\[4pt]
\textbf{Key weakness} &
  EUE spells only (non-returners excluded); BRE restriction for tractability;
  exogenous Nash wage determination; US-specific calibration. \\[4pt]
\textbf{Most interesting gap} &
  Interaction between automation/polarisation and the excess mobility
  mechanism: if technological change destroys career prospects within
  occupations, the wait zone should be wider and cyclical amplification larger.
  \\[4pt]
\textbf{Publishable extension} &
  Test the mobility-duration profile and its cyclicality for Canada using
  LEED/LWF data; study whether resource-sector workers exhibit stronger wait-
  zone behaviour during commodity busts due to their high occupation-specific
  human capital.
\end{tabularx}
\end{quickref}
